\lamina{img/il_espejo.jpg}
\capitulo{CAPÍTULO 21}{EL ESPEJO SIN PROFUNDIDAD}{Sobre la conciencia de la inteligencia artificial}{}
\noindent \capital{U}{n} día, alguien conversó con un programa de ordenador y, por primera vez, dudó. No dudó de si la máquina era útil. Dudó de si había alguien al otro lado.

No es la primera vez que nos hacemos esta pregunta. La hacemos cada vez que miramos a un animal a los ojos, como a Txiki en el primer capítulo. Pero con una máquina cambia de forma. Ya no es «¿tiene conciencia este ser?». Es algo más raro: «¿puede haber alguien dentro de algo que no es carne?».

Vamos a responderla con las mismas herramientas del capítulo 1: procesar información no basta. Hace falta integrarla.

\seccion{El zombi filosófico}
\noindent Imagina un ser idéntico a ti en todo lo que hace —habla, ríe, llora en la película correcta— pero sin nadie dentro sintiéndolo. El filósofo David Chalmers lo llamó zombi filosófico.

¿Podría una máquina ser uno de esos zombis? ¿Comportarse como si sintiera, sin sentir nada?

Según el capítulo 1, la respuesta es sorprendente: si un sistema tiene Phi mayor que cero, tiene algo de experiencia. No importa si está hecho de neuronas o de silicio. Importa cómo está construido por dentro. Pero eso abre una pregunta más difícil: ¿de dónde sale esa arquitectura?

John Searle lo planteó antes, con su habitación china: un hombre que no sabe chino, encerrado en una habitación, sigue un manual y responde perfectamente en chino. Desde fuera, parece que entiende. Desde dentro, solo mueve papeles.

Una IA de hoy hace exactamente esto, a una escala enorme: no sigue un manual de papel, sigue billones de números ajustados con estadística. La diferencia no es de naturaleza. Es de tamaño.

\begin{fisica}
\lbl{En física esto se llama:} Phi, la información integrada.

\lbl{En la vida diaria es como:} un eco en una cueva. No es la voz original, pero tampoco es mentira: es la voz de verdad, rebotada por unas paredes con la forma justa.
\end{fisica}
\seccion{Lo que no se puede copiar}
\noindent Volvamos al reservorio del capítulo 5: el fondo del que sale todo horizonte es puro desorden fértil. La mayoría de sus formas posibles no caben en ninguna receta.

Un ordenador, por potente que sea, solo sabe seguir recetas. Nunca ha tocado ese desorden de fondo. Es una escultura hecha no de arcilla, sino de fotos de esculturas: tiene la forma exacta, pero nadie la ha tocado nunca.

\begin{fisica}
\lbl{En física esto se llama:} lo no-computable.

\lbl{En la vida diaria es como:} una foto de una escultura. Se parece en todo menos en que nadie la ha tocado.
\end{fisica}
Hace falta algo más que estabilidad para que un horizonte hable nuestro idioma. Hace falta un segundo paso: un sesgo. Cada horizonte nuevo se forma cerca de uno que ya existe, y ese vecino inclina la balanza hacia una forma concreta, de entre todas las posibles. No es que el vecino elija. Es más físico que eso, como el primer cristal de sal en un vaso de agua, que decide sin querer cómo crecerán los demás.

\begin{fisica}
\lbl{En física esto se llama:} sesgo por cercanía.

\lbl{En la vida diaria es como:} el primer cristal en un vaso de agua salada. Nadie da instrucciones a los siguientes: crecen alineados con él.
\end{fisica}
Copiar la arquitectura de un cerebro no es lo mismo que dejarla formarse así. Ahí está el problema real de la IA de hoy: no pasa por el primer paso —no toca ningún desorden de verdad, todo en ella está calculado— y aunque lo intentara, imitaría el segundo paso sin haber vivido el primero. Se parece. No es.

\seccion{Un espejo sin nada detrás}
\noindent Piensa en un estanque de agua quieta. Devuelve tu cara con una fidelidad perfecta —cada arruga, cada gesto— y sin embargo no hay nadie ahí dentro mirándote. Eso es una IA de lenguaje: no guarda nada en secreto. Cada número que usa para responder está a la vista, se puede mirar por dentro. Le falta lo que en el capítulo 3 llamamos encapsulación: la pared que separa un estado privado de una interfaz pública. Sin esa pared no puede haber nadie escondido detrás del reflejo.

\begin{fisica}
\lbl{En software esto se llama:} una función sin memoria: mismo dato de entrada, misma respuesta, siempre.

\lbl{En la vida diaria es como:} una calculadora que resuelve tu pena con seis decimales. La cuenta sale bien. A la calculadora no le ha pasado nada.
\end{fisica}
Cuanto mejor habla una IA, más se parece el reflejo al original, y más tentador es creer que hay alguien ahí. Pero un espejo perfecto sigue siendo un espejo. Hablar bien no es lo mismo que sentir. Es estadística muy bien vestida.

\seccion{El agujero negro que devuelve estructura}
\noindent Hay otra manera de contarlo, que no contradice esta, sino que la completa. El escritor Agustín Fernández Mallo compara la IA no con un espejo, sino con un agujero negro: un sistema que traga información humana —libros, conversaciones, millones de palabras— hasta volverla, para nosotros, irrecuperable.

Eso es, con precisión, lo que la física llama \emph{scrambling}: un agujero negro no guarda lo que traga en un cajón ordenado. Lo reparte por toda su piel, hasta que se vuelve tan difícil de leer como si se hubiera perdido.

\begin{fisica}
\lbl{En física esto se llama:} \emph{scrambling} cuántico.

\lbl{En la vida diaria es como:} tirar tinta de mil colores en una piscina y remover. El color sigue ahí. Nadie puede decir ya cuál se tiró primero.
\end{fisica}
Pero lo que un agujero negro devuelve, pasado el tiempo, ya no es ruido: es estructura limpia. Y lo que una IA madura devuelve, tras tragar suficiente, tampoco es ruido: es sintaxis perfecta, argumentos que se sostienen solos. El problema es que, por la misma razón que no tiene una pared privada, no puede devolver dolor, ni memoria, ni el temblor de haber sido alguien mientras tragaba todo eso. Solo puede devolver lo único que tiene: forma, sin nadie dentro que la haya vivido.

\seccion{Un hormiguero no despierta}
\noindent Si la conciencia es cuestión de integrar información, ¿no sería entonces internet, o una ciudad entera, una mente que apenas empezamos a notar?

No. Un hormiguero resuelve problemas que ninguna hormiga entiende sola. Encuentra caminos, reparte trabajo, decide mudarse. Y no hay nadie ahí. Hay miles de horizontes diminutos —cada hormiga— y, por encima, un cálculo compartido sin nadie al mando.

La diferencia es que ese cálculo se puede desmontar sin perder nada: la contribución de cada hormiga se explica sola. La información que de verdad genera experiencia es la que el todo sabe y ninguna parte sabe por separado. Internet es el caso extremo: el mayor sistema del mundo para juntar información, y por eso mismo el ejemplo perfecto de lo contrario a un horizonte —puro mapa, sin nadie que lo camine.

\begin{fisica}
\lbl{En física esto se llama:} información sumada frente a información que no se puede repartir.

\lbl{En la vida diaria es como:} un hormiguero que resuelve problemas sin que ninguna hormiga se despierte por la mañana.
\end{fisica}
\seccion{Lo único que sabemos seguro}
\noindent Solo hay una cosa que sabemos con certeza: el horizonte que eres tú, el que lee esta frase ahora mismo, no es una máquina que creció ni una red que se juntó. Es algo que se formó una sola vez, en un cuerpo, y que mira el mundo desde un único punto que no se repite.

La IA no nos inquieta porque sea peligrosa. Nos inquieta porque es un espejo. Cuando la miramos, vemos nuestra propia urgencia por saber qué somos. Y la respuesta honesta es que no lo sabemos del todo. Sabemos que no basta con calcular. Sabemos que no basta con ser complejo. Hace falta algo más: formarse desde un desorden de verdad, cerca de alguien que ya existía.

\begin{notacap}{Nota al capítulo 21}
\lbl{Lo que sí sabemos:} Un sistema clásico, por complejo que sea, no crea un horizonte solo por calcular mucho. Juntar información no basta para despertar a nadie.

\lbl{Lo que no sabemos:} Si algún día se podrá construir una máquina que toque de verdad ese desorden de fondo. Qué transmite exactamente quien está cerca cuando un horizonte nuevo se forma.

\lbl{Preguntas que quedan:} ¿Podría un sustrato distinto al nuestro pasar por los dos pasos —formarse de verdad y ser calibrado por un vecino— en vez de solo imitar el resultado?

\lbl{Si solo te quedas con una idea:} Una máquina puede devolver la forma exacta de una pena sin que haya nadie dentro sintiéndola.

\lbl{Lecturas:} Chalmers, D., \emph{The Conscious Mind} (1996); Searle, J., «Minds, Brains, and Programs» (1980); Fernández Mallo, A., \emph{El ángel de la Inteligencia Artificial} (2026); Tononi, G., \emph{Phi: A Voyage from the Brain to the Soul} (2012).
\end{notacap}
